\documentclass[11pt,twoside,final,openright]{xenstyle}
\usepackage{a4,graphicx,setspace,times}
\setstretch{1.15}
%
\usepackage{parskip,fancyheadings}

\setcounter{secnumdepth}{10}
\setcounter{tocdepth}{2}

%\usepackage{sabon}
\def\textos#1{#1}

%
% From PRB: Use FancyPlain style -- set up to produce draft
% information and disable rules and section headings in headers and
% footers.
%

% \filename

\def\filename{thesis.tex}
\let\savedinclude\include
\def\include#1{\def\filename{#1.tex}\savedinclude{#1}}

% \prtinttime

\newcounter{hours}\newcounter{mins}
\newcommand{\printtime}{%
  \setcounter{hours}{\time/60}%
  \setcounter{mins}{\time-\value{hours}*60}%
  \ifthenelse{\value{hours}<10}{0}{}\thehours:%
  \ifthenelse{\value{mins}<10}{0}{}\themins}


% Final tweaks to the layout -- PRB/SMH/TLH

\newcommand{\verylongpage}{\enlargethispage{2.5\baselineskip}}
\newcommand{\longpage}{\enlargethispage{2\baselineskip}}
\newcommand{\shortpage}{\enlargethispage{-2\baselineskip}}
\newcommand{\quitelongpage}{\enlargethispage{1\baselineskip}}
\newcommand{\quiteshortpage}{\enlargethispage{-1\baselineskip}}

% Use instead of / to allow hyphenation -- PRB/SMH

\def\slash{/\penalty 0\hskip 0pt\relax}

%
% Header/Footer
%

\pagestyle{fancyplain}
\headrulewidth=0pt
\footrulewidth=0pt
\plainheadrulewidth=0pt
\plainfootrulewidth=0pt

\lhead[]{}
\chead[]{}
\rhead[]{}

\lfoot[\sl\thepage]{}
\cfoot[]{}
\rfoot[]{\sl\thepage}

%
% Captions (from RF: make captions hang properly)
%

\makeatletter
\long\def\@makecaption#1#2{%
   \vskip 10\p@
   \setbox\@tempboxa\hbox{\small #1: #2}%
   \ifdim \wd\@tempboxa >\hsize
       \@hangfrom{\small #1: }{\small #2}\par
     \else
       \hbox to\hsize{\hfil\box\@tempboxa\hfil}%
   \fi}
\makeatother

\makeatletter

%
% Keep mpars on the same side (we use them for hanging section
% numbers).
%

\@mparswitchfalse
\marginparsep=1ex

%
% Abandon ugly vertical justification
%

\raggedbottom

%
% Poor-mans small caps
%

%\font \psc = ptmr8t scaled 800
%\renewcommand{\scdefault}{psc}

\DeclareOldFontCommand{\sc}{\normalfont\scshape}{\@nomath\sc}

%
% Let footnotes fall right to the bottom of the page
%

\renewcommand\footnoterule{\vfill\hrule\vspace{6pt}}

%%%%%%%%%%%%%%%%%%%%%%%%%%%%%%%%%%%%%%%%%%%%%%%%%%%%%%%%%%%%%%%%%%%%%%
%
% Headings

%
% Fonts
%

% \font\sectionfont = phvr8t scaled 1400
% \font\subsectionfont = phvr8t scaled 1200
% \font\subsubsectionfont = phvr8t scaled 1000

%
% Chapter/Section names
%

\def\chaptermark#1{%
        \markboth {{\ifnum \c@secnumdepth >\m@ne
            \@chapapp\ \thechapter. \ \fi
            #1}}{}
        }

\def\sectionmark#1{%
        \markright {{\ifnum \c@secnumdepth >\z@
            \thesection. \ \fi
            #1}}
        }

%
% Chapter heading components
%

\def\@chapter[#1]#2{\ifnum \c@secnumdepth >\m@ne
                         \refstepcounter{chapter}%
                         \typeout{\@chapapp\space\thechapter.}%
                         \addcontentsline{toc}{chapter}%
                                   {\protect\numberline{\thechapter}#1}%
                    \else
                      \addcontentsline{toc}{chapter}{#1}
                    \fi
                    \chaptermark{#1}%
                    \addtocontents{lof}{\protect\addvspace{10\p@}}%
                    \addtocontents{lot}{\protect\addvspace{10\p@}}%
                    \if@twocolumn
                      \@topnewpage[\@makechapterhead{#2}]%
                    \else
                      \@makechapterhead{#2}%
                      \@afterheading
                    \fi
        }

\def\@makechapterhead#1{%
  \vspace*{50\p@}%
  {\parindent \z@ \raggedright \reset@font
    \ifnum \c@secnumdepth >\m@ne
         \LARGE \@chapapp{} \thechapter
         \par
         \vskip 10\p@
       \fi
    \Huge #1\par
    \nobreak
    \vskip 20\p@
  }}

\def\@schapter#1{\if@twocolumn
                   \@topnewpage[\@makeschapterhead{#1}]%
                 \else
                   \@makeschapterhead{#1}%
                   \@afterheading
                 \fi}

\def\@makeschapterhead#1{%
  \vspace*{50\p@}%
  {\parindent \z@ \raggedright
    \reset@font
    \Huge #1\par
    \nobreak
    \vskip 40\p@
  }}

%
% *Section heading components
%

\def\@startsection#1#2#3#4#5#6{%
  \if@noskipsec \leavevmode \fi
  \par
  \@tempskipa #4\relax
  \@afterindenttrue
  \ifdim \@tempskipa <\z@
    \@tempskipa -\@tempskipa \@afterindentfalse
  \fi
  \if@nobreak
    \everypar{}%
  \else
    \addpenalty\@secpenalty\addvspace\@tempskipa
  \fi
  \@ifstar
    {\@ssect{#3}{#4}{#5}{#6}}%
    {\@dblarg{\@sect{#1}{#2}{#3}{#4}{#5}{#6}}}}

\def\@sect#1#2#3#4#5#6[#7]#8{%
  \ifnum #2>\c@secnumdepth
    \let\@svsec\@empty
  \else
    \refstepcounter{#1}%
    \protected@edef\@svsec{\@seccntformat{#1}\relax}%
  \fi
  \@tempskipa #5\relax
  \ifdim \@tempskipa>\z@
    \begingroup
      #6{\interlinepenalty \@M #8\@@par}%
      \reversemarginpar{\marginpar{{\rightline{\@svsec}}}}
    \endgroup
    \csname #1mark\endcsname{#7}%
    \addcontentsline{toc}{#1}{%
      \ifnum #2>\c@secnumdepth \else
        \protect\numberline{\csname the#1\endcsname}%
      \fi
      #7}%
  \else
    \def\@svsechd{%
      #6{\hskip #3\relax
      \@svsec #8}%
      \csname #1mark\endcsname{#7}%
      \addcontentsline{toc}{#1}{%
        \ifnum #2>\c@secnumdepth \else
          \protect\numberline{\csname the#1\endcsname}%
        \fi
        #7}}%
  \fi
  \@xsect{#5}
  }

%
% Table of contents
%

\renewcommand\contentsname{Table of contents}
\renewcommand\tableofcontents{%
    \if@twocolumn
      \@restonecoltrue\onecolumn
    \else
      \@restonecolfalse
    \fi
    \chapter*{\contentsname
        \@mkboth{\contentsname}{}}%
    \@starttoc{toc}%
    \if@restonecol\twocolumn\fi
    }

%
% List of Figures
%

\renewcommand\listfigurename{List of figures}
\renewcommand\listoffigures{%
    \if@twocolumn
      \@restonecoltrue\onecolumn
    \else
      \@restonecolfalse
    \fi
    \chapter*{\listfigurename
      \@mkboth{\listfigurename}{}}%
    \@starttoc{lof}%
    \if@restonecol\twocolumn\fi
    }

%
% List of Tables
%

\renewcommand\listtablename{List of tables}
\renewcommand\listoftables{%
    \if@twocolumn
      \@restonecoltrue\onecolumn
    \else
      \@restonecolfalse
    \fi
    \chapter*{\listtablename
      \@mkboth{\listtablename}{}}%
    \@starttoc{lot}%
    \if@restonecol\twocolumn\fi
    }

%
% Glossary
%

\newcommand\gloname{Glossary}
\newcommand\gloss{%
    \if@twocolumn
      \@restonecoltrue\onecolumn
    \else
      \@restonecolfalse
    \fi
    \chapter*{\gloname
      \@mkboth{\gloname}{}}%
    \chapter{Glossary of Terms}

\begin{description}

\item[Atropos] One of the CPU schedulers provided by Xen.  Atropos
  provides domains with absolute shares of the CPU, with timeliness
  guarantees and a mechanism for sharing out `slack time'.

\item[BVT] The BVT scheduler is used to give proportional fair shares
  of the CPU to domains.

\item[Exokernel] A minimal piece of privileged code, similar to a {\bf
    microkernel} but providing a more `hardware-like' interface to the
  tasks it manages.  This is similar to a paravirtualising VMM like
  {\bf Xen} but was designed as a new operating system structure,
  rather than specifically to run multiple conventional OSs.

\item[Domain] A domain is the execution context that contains a
  running {\bf virtual machine}.  The relationship between virtual
  machines and domains on Xen is similar to that between programs and
  processes in an operating system: a virtual machine is a persistent
  entity that resides on disk (somewhat like a program).  When it is
  loaded for execution, it runs in a domain.  Each domain has a {\bf
    domain ID}.

\item[Domain 0] The first domain to be started on a Xen machine.
  Domain 0 is responsible for managing the system.

\item[Domain ID] A unique identifier for a {\bf domain}, analogous to
  a process ID in an operating system.

\item[Full virtualisation] An approach to virtualisation which
  requires no modifications to the hosted operating system, providing
  the illusion of a complete system of real hardware devices.

\item[Hypervisor] An alternative term for {\bf VMM}, used because it
  means `beyond supervisor', since it is responsible for managing
  multiple `supervisor' kernels.

\item[Live migration] A technique for moving a running virtual machine
  to another physical host, without stopping it or the services
  running on it.

\item[Microkernel] A small base of code running at the highest
  hardware privilege level.  A microkernel is responsible for sharing
  CPU and memory (and sometimes other devices) between less privileged
  tasks running on the system.  This is similar to a VMM, particularly
  a {\bf paravirtualising} VMM but typically addressing a different
  problem space and providing different kind of interface.

\item[NetBSD/Xen] A port of NetBSD to the Xen architecture.

\item[Paravirtualisation] An approach to virtualisation which requires
  modifications to the operating system in order to run in a virtual
  machine.  Xen uses paravirtualisation but preserves binary
  compatibility for user space applications.

\item[Shadow pagetables] A technique for hiding the layout of machine
  memory from a virtual machine's operating system.  Used in some {\bf
  VMMs} to provide the illusion of contiguous physical memory, in
  Xen this is used during {\bf live migration}.

\item[Virtual Block Device] Persistant storage available to a virtual
  machine, providing the abstraction of an actual block storage device.
  {\bf VBD}s may be actual block devices, filesystem images, or
  remote/network storage.

\item[Virtual Machine] The environment in which a hosted operating
  system runs, providing the abstraction of a dedicated machine.  A
  virtual machine may be identical to the underlying hardware (as in
  {\bf full virtualisation}, or it may differ, as in {\bf
  paravirtualisation}).

\item[VMM] Virtual Machine Monitor - the software that allows multiple
  virtual machines to be multiplexed on a single physical machine.

\item[Xen] Xen is a paravirtualising virtual machine monitor,
  developed primarily by the Systems Research Group at the University
  of Cambridge Computer Laboratory.

\item[XenLinux] Official name for the port of the Linux kernel that
  runs on Xen.

\end{description}

    \vfill
    \if@restonecol\twocolumn\fi
    }

%
% Bibliography
%

\renewenvironment{thebibliography}[1]
     {\chapter*{\bibname
        \@mkboth{\bibname}{}}%
      \list{\@biblabel{\arabic{enumiv}}}%
           {\settowidth\labelwidth{\@biblabel{#1}}%
            \leftmargin\labelwidth
            \advance\leftmargin\labelsep
%            \if@openbib
%              \advance\leftmargin\bibindent
%              \itemindent -\bibindent
%              \listparindent \itemindent
%              \parsep \z@
%            \fi
            \usecounter{enumiv}%
            \let\p@enumiv\@empty
            \renewcommand\theenumiv{\arabic{enumiv}}}%
%      \if@openbib
%        \renewcommand\newblock{\par}
%      \else
        \renewcommand\newblock{\hskip .11em \@plus.33em \@minus.07em}%
%      \fi
      \sloppy\clubpenalty4000\widowpenalty4000%
      \sfcode`\.=\@m}
     {\def\@noitemerr
       {\@latex@warning{Empty `thebibliography' environment}}%
      \endlist}

%
% Chapter
%

\renewcommand\chapter{%
    \if@openright\cleardoublepage\else\clearpage\fi
    \thispagestyle{plain}%
    \global\@topnum\z@
    \@afterindentfalse
    \secdef\@chapter\@schapter
    }

%
% *Sections
%

% \z@ -- zero?

% \@startsection{name}{depth}
%         {indent}
%         {beforeskip \@plus foo \@minus bar} 
%         foo: ... ; bar: jitter back up page?
%         {afterskip}
%         {style}

\renewcommand\section{%
    \@startsection{section}{1}
        {\z@}%
        {-5ex \@plus 5ex \@minus -.5ex}%
        {2.5ex  \@minus .5ex}%
        {\reset@font\Large}%
    }

\renewcommand\subsection{%
    \@startsection{subsection}{2}
        {\z@}%
        {-3ex  \@plus 3ex \@minus -.5ex}%
        {1.5ex \@minus .5ex}%
        {\reset@font\large}%
    }

\renewcommand\subsubsection{%
    \@startsection{subsubsection}{3}
        {\z@}%
        {-1ex \@plus 1ex \@minus -.5ex}%
        {1ex  \@minus .5ex}%
        {\reset@font\normalsize\sl}%
    }

\makeatother
\endinput



\begin{document}

% TITLE PAGE
\pagestyle{empty}
\begin{center}
\vspace*{\fill}
\includegraphics{figs/xenlogo.eps}
\vfill
\vfill
\vfill
\begin{tabular}{l}
{\Huge \bf Users' manual} \\[4mm]
{\huge Xen v2.0 for x86} \\[80mm]

{\Large Xen is Copyright (c) 2004, The Xen Team} \\[3mm]
{\Large University of Cambridge, UK} \\[20mm]
{\large Last updated on 26th October, 2004}
\end{tabular}
\vfill
\end{center}
\cleardoublepage

% TABLE OF CONTENTS
\pagestyle{plain}
\pagenumbering{roman}
{ \parskip 0pt plus 1pt
  \tableofcontents }
\cleardoublepage

% PREPARE FOR MAIN TEXT
\pagenumbering{arabic}
\raggedbottom
\widowpenalty=10000
\clubpenalty=10000
\parindent=0pt
\parskip=5pt
\renewcommand{\topfraction}{.8}
\renewcommand{\bottomfraction}{.8}
\renewcommand{\textfraction}{.2}
\renewcommand{\floatpagefraction}{.8}
\setstretch{1.15}

\newcommand{\path}[1]{{\tt #1}}

\part{Introduction and Tutorial}
\chapter{Introduction}

{\bf
DISCLAIMER: This documentation is currently under active development
and as such there may be mistakes and omissions --- watch out for
these and please report any you find to the developer's mailing list.
Contributions of material, suggestions and corrections are welcome.
}

\vspace{5mm}

Xen is a { \em paravirtualising } virtual machine monitor (VMM), or
`hypervisor', for the x86 processor architecture.  Xen can securely
execute multiple virtual machines on a single physical system with
close-to-native performance.  The virtual machine technology
facilitates enterprise-grade functionality, including:

\begin{itemize}
\item Virtual machines with performance nearly identical to native
  hardware.
\item Live migration of running virtual machines.
\item Excellent hardware support (supports most Linux device drivers).
\item Suspend to disk / resume from disk of running virtual machines.
\item Transparent copy-on-write disks for simple sharing across VMs.
\item Sandboxed, restartable device drivers.
\end{itemize}

Paravirtualisation allows very high performance virtual machine
technology, even on architectures like x86 that are traditionally
hard to virtualise.
The drawback of this approach is that it requires operating systems to
be {\em ported} to run on Xen.  This process is similar to a port of
an operating system to a new hardware platform, although the process
is simplified because the paravirtual machine architecture is very
similar to the underlying native hardware. Although operating system
kernels must explicitly support Xen, a key feature is that user space
applications and libraries {\em do not} require modification.

Xen support is available for increasingly many operating systems:
right now, Linux 2.4, Linux 2.6 and NetBSD are available for Xen 2.0.
We expect that Xen support will ultimately be integrated into the
official releases of Linux, NetBSD, FreeBSD and Dragonfly BSD. 
Other OS ports, including Plan 9, are in progress.

%Even running multiple copies of Linux can be very useful, providing a
%means of containing faults to one OS image, providing performance
%isolation between the various OS instances and trying out multiple
%distros.

Possible usage scenarios for Xen include:
\begin{description}
\item [Kernel development.] Test and debug kernel modifications in a
      sandboxed virtual machine --- no need for a separate test
      machine.
\item [Multiple OS configurations.] Run multiple operating systems
      simultaneously, for instance for compatibility or QA purposes.
\item [Server consolidation.] Move multiple servers onto one box,
      provided performance and fault isolation at virtual machine
      boundaries.
\item [Cluster computing.] Improve manageability and efficiency by
      running services in virtual machines, isolated from
      machine-specifics and load balance using live migration.
\item [High availability computing.] Run device drivers in sandboxed
      domains for increased robustness.
\item [Hardware support for custom OSes.] Export drivers from a
      mainstream OS (e.g. Linux) with good hardware support
      to your custom OS, avoiding the need for you to port existing
      drivers to achieve good hardware support.
\end{description}

\section{Structure of a Xen-Based System}

A Xen system has multiple layers, the lowest and most privileged of
which is Xen itself. Xen in turn may host multiple {\em guest}
operating systems, each of which is executed within a secure virtual
machine (in Xen terminology, a {\em domain}). Domains are scheduled by
Xen to make effective use of the available physical CPUs.  Each guest
OS manages its own applications, which includes responsibility for
scheduling each application within the time allotted to the VM by Xen.

The first domain, {\em Domain 0}, is created automatically when the
system boots and has special management privileges. Domain 0 builds
other domains and manages their virtual devices. It also performs
suspend, resume and migration of virtual machines. Where one is
required, the X server is also run in domain 0.

Within Domain 0, a process called `Xend' runs to manage the system.
Xend is responsible for managing virtual machines and providing access
to their consoles.  Commands are issued to Xend over an HTTP
interface, either from a command-line tool or from a web browser.

\section{Hardware Support}

Xen currently runs only on the x86 architecture (however, ports to other
architectures, including x86/64 and IA64, are in progress).
Xen requires a `P6' or newer processor (e.g. Pentium Pro, Celeron,
Pentium II, Pentium III, Pentium IV, Xeon, AMD Athlon, AMD Duron).
Multiprocessor machines are supported, and we also have basic support
for HyperThreading (SMT), although this remains a topic for ongoing
research.  As mentioned above, a port specifically for x86/64 is in
progress, although Xen already runs on such systems in 32-bit legacy
mode.

Xen can currently use up to 4GB of memory.  It is possible for x86
machines to address up to 64GB of physical memory but there are no
plans to support these systems.  The x86\_64 port is the planned route
to supporting more than 4GB of memory.

Xen offloads most of the hardware support issues to the guest OS
running in Domain 0.  Xen itself only contains code to detect and
start additional processors, set up interrupt routing and perform PCI
bus enumeration.  Device drivers run within a privileged guest OS
rather than within Xen itself, meaning that we should be
compatible with the majority of device hardware supported by Linux.
The default XenLinux build contains support for relatively modern
server-class network and disk hardware, but you can add support for
other hardware by configuring your XenLinux kernel in the normal way
(e.g. \verb_# make ARCH=xen menuconfig_).

\section{History}


Xen was originally developed by the Systems Research Group at the
University of Cambridge Computer Laboratory as part of the XenoServers
project, funded by UK-EPSRC.
XenoServers aim to provide a `public infrastructure for
global distributed computing', and Xen plays a key part in that,
allowing us to efficiently partition a single machine to enable
multiple independent clients to run their operating systems and
applications in an environment providing protection, resource
isolation and accounting.  The project web page contains further
information along with pointers to papers and technical reports:
{\tt http://www.cl.cam.ac.uk/xeno}

Xen has since grown into a project in its own right, enabling us to
investigate interesting research issues regarding the best techniques
for virtualising resources such as the CPU, memory, disk and network.
The project has been bolstered by support from Intel Research
Cambridge, and HP Labs, who are now working closely with us.

Xen was first described in the 2003 paper at SOSP\footnote{\tt
http://www.cl.cam.ac.uk/netos/papers/2003-xensosp.pdf}, and the
first public release (1.0) was made in October 2003.  Since
then, Xen has been extensively developed and is now used in production
scenarios on multiple sites.

Xen 2.0 feature greatly enhanced hardware support, configuration
flexibility, usability and a larger complement of supported operating
systems. We think that Xen has the potential to become {\em the}
definitive open source virtualisation solution and will work to
conclusively achieve that position.

\chapter{Installation}

The Xen distribution includes three main components:  Xen itself,
utilities to convert a standard Linux tree to run on Xen and the
userspace tools required to operate a Xen-based system.

This manual describes how to install the Xen 2.0 distribution from
source.  Alternatively, there may be packages available for your
operating system distribution.

\section{Prerequisites}
\label{sec:prerequisites}
\begin{itemize}
\item i686-class CPU or newer
\item A working installation of your favourite Linux distribution.
\item A working installation of the GRUB bootloader.
\item An installation of Twisted v1.3 or above (see {\tt
http://www.twistedmatrix.com}).  There may be a package available for
your distribution; alternatively it can be installed by running {\tt \#
make install-twisted} in the root of the Xen source tree.
\item Python logging package (see {\tt http://www.red-dove.com/})
\item The Linux bridge control tools (see {\tt
http://bridge.sourceforge.net}).  There may be packages of these tools
available for your distribution.
\item Linux IP Routing Tools
\item make
\item gcc
\item libcurl
\item zlib-dev
\item python-dev
\item python2.3-pycurl
\item python2.3-twisted
\end{itemize}

\section{Install Bitkeeper (Optional)}

To fetch a local copy, first download the BitKeeper tools.
Download instructions must be obtained by filling out the provided
form at: \\ {\tt
http://www.bitmover.com/cgi-bin/download.cgi }

The BitKeeper install program is designed to be run with X.  If X is
not available, you can specify the install directory on the command
line.

\section{Download the Xen source code}

\subsection{Using Bitkeeper}

The public master BK repository for the 2.0 release lives at: \\
{\tt bk://xen.bkbits.net/xen-2.0.bk}.  You can use Bitkeeper to
download it and keep it updated with the latest features and fixes.

Change to the directory in which you want to put the source code, then
run:
\begin{verbatim}
# bk clone bk://xen.bkbits.net/xen-2.0.bk
\end{verbatim}

Under your current directory, a new directory named `xen-2.0.bk' has
been created, which contains all the source code for the Xen
hypervisor and the Xen tools.  The directory also contains `sparse' OS
source trees, containing only the files that require changes to allow
the OS to run on Xen.

Once you have cloned the repository, you can update to the newest
changes to the repository by running:
\begin{verbatim}
# cd xen-2.0.bk # to change into the local repository
# bk pull       # to update the repository
\end{verbatim}

\subsection{Without Bitkeeper}

The Xen source tree is also available in gzipped tarball form from the
Xen downloads page:\\
{\tt http://www.cl.cam.ac.uk/Research/SRG/netos/xen/downloads.html}.
Prebuilt tarballs are also available from this page but are very large.

\section{The distribution}

The Xen source code repository is structured as follows:

\begin{description}
\item[\path{tools/}] Xen node controller daemon (Xend), command line tools, 
  control libraries
\item[\path{xen/}] The Xen hypervisor itself.
\item[\path{linux-2.4.27-xen-sparse/}] Xen support for Linux 2.4
\item[\path{linux-2.6.9-xen-sparse/}] Xen support for Linux 2.6
\item[\path{linux-2.6.9-patches/}] Experimental patches for Linux 2.6
\item[\path{netbsd-2.0-xen-sparse/}] Xen support for NetBSD 2.0
\item[\path{docs/}] various documentation files for users and developers
\item[\path{extras/}] currently this contains the Mini OS, aimed at developers
\end{description}

\section{Build and install}

The Xen makefile includes a target `world' that will do the
following:

\begin{itemize}
\item Build Xen
\item Build the control tools, including Xend
\item Download (if necessary) and unpack the Linux 2.6 source code,
      and patch it for use with Xen
\item Build a Linux kernel to use in domain 0 and a smaller
      unprivileged kernel, which can optionally be used for
      unprivileged virtual machines.
\end{itemize}

Inspect the Makefile if you want to see what goes on during a build.
Building Xen and the tools is straightforward, but XenLinux is more
complicated.  The makefile needs a `pristine' Linux kernel tree which
it will then add the Xen architecture files to.  You can tell the
makefile the location of the appropriate Linux compressed tar file by
setting the LINUX\_SRC environment variable, e.g. \\
\verb!# LINUX_SRC=/tmp/linux-2.6.8.1.tar.bz2 make world! \\ or by
placing the tar file somewhere in the search path of {\tt
LINUX\_SRC\_PATH} which defaults to `{\tt .:..}'.  If the makefile
can't find a suitable kernel tar file it attempts to download it from
kernel.org (this won't work if you're behind a firewall).

After untaring the pristine kernel tree, the makefile uses the {\tt
mkbuildtree} script to add the Xen patches to the kernel.  It then
builds two different XenLinux images, one with a `-xen0' extension
which contains hardware device drivers and drivers for Xen's virtual
devices, and one with a `-xenU' extension that just contains the
virtual ones.

The procedure is similar to build the Linux 2.4 port: \\
\verb!# LINUX_SRC=/path/to/linux2.4/source make linux24!

The NetBSD port can be built using: \\ \verb!# make netbsd! \\ The
NetBSD port is built using a snapshot of the netbsd-2-0 cvs branch.
The snapshot is downloaded as part of the build process, if it is not
yet present in the {\tt NETBSD\_SRC\_PATH} search path.  The build
process also downloads a toolchain which includes all the tools
necessary to build the NetBSD kernel under Linux.

If you have an SMP machine you may wish to give the {\tt '-j4'}
argument to make to get a parallel build.

If you have an existing Linux kernel configuration that you would like
to use for domain 0, you should copy it to
install/boot/config-2.6.8.1-xen0.  During the first build, you may be
asked about some Xen-specific options.  We advise accepting the
defaults for these options.

\framebox{\parbox{5in}{
{\bf Distro specific:} \\
{\it Gentoo} --- if not using udev (most installations, currently), you'll need
to enable devfs and devfs mount at boot time in the xen0 config.
}}

The files produced by the build process are stored under the
\path{install/} directory.  To install them in their default
locations, do: \\
\verb_# make install_

Alternatively, users with special installation requirements may wish
to install them manually by copying the files to their appropriate
destinations.

Files in \path{install/boot/} include:
\begin{itemize}
\item \path{install/boot/xen.gz} The Xen 'kernel'
\item \path{install/boot/vmlinuz-2.6.8.1-xen0}  Domain 0 XenLinux kernel
\item \path{install/boot/vmlinuz-2.6.8.1-xenU}  Unprivileged XenLinux kernel
\end{itemize}

The difference between the two Linux kernels that are built is due to
the configuration file used for each.  The "U" suffixed unprivileged
version doesn't contain any of the physical hardware device drivers
--- it is 30\% smaller and hence may be preferred for your
non-privileged domains.  The `0' suffixed privileged version can be
used to boot the system, as well as in driver domains and unprivileged
domains.

The \path{install/boot} directory will also contain the config files
used for building the XenLinux kernels, and also versions of Xen and
XenLinux kernels that contain debug symbols (\path{xen-syms} and
\path{vmlinux-syms-2.6.8.1-xen0}) which are essential for interpreting crash
dumps.  Retain these files as the developers may wish to see them if
you post on the mailing list.

\section{Configuration}

\subsection{GRUB Configuration}

An entry should be added to \path{grub.conf} (often found under
\path{/boot/} or \path{/boot/grub/}) to allow Xen / XenLinux to boot.
This file is sometimes called \path{menu.lst}, depending on your
distribution.  The entry should look something like the following:

\begin{verbatim}
title Xen 2.0 / XenLinux 2.6.8.1
  kernel /boot/xen.gz dom0_mem=131072 com1=115200,8n1
  module /boot/vmlinuz-2.6.8.1-xen0 root=/dev/sda4 ro console=tty0 console=ttyS0
\end{verbatim}

The first line of the configuration (kernel...) tells GRUB where to
find Xen itself and what boot parameters should be passed to it (in
this case, setting domain 0's memory allocation and the settings for
the serial port).

The second line of the configuration describes the location of the
XenLinux kernel that Xen should start and the parameters that should
be passed to it (these are standard Linux parameters, identifying the
root device and specifying it be initially mounted read only and
instructing that console output be sent both to the screen and to the
serial port).

If you want to use an initrd, just add another {\tt module} line to
the configuration, as usual:
\begin{verbatim}
  module /boot/my_initrd.gz
\end{verbatim}

As always when installing a new kernel, it is recommended that you do
not remove the original contents of \path{menu.lst} --- you may want
to boot up with your old Linux kernel in future, particularly if you
have problems.

\framebox{\parbox{5in}{
{\bf Distro specific:} \\
{\it SuSE} --- Omit the {\tt ro} option from the XenLinux kernel
command line, since the partition won't be remounted rw during boot.
}}

\subsection{Serial Console}

In order to configure serial console output, it is necessary to add a
line into \path{/etc/inittab}.  The XenLinux console driver is
designed to make this procedure the same as configuring a normal
serial console.  Add the line:

{\tt c:2345:respawn:/sbin/mingetty ttyS0}

\subsection{TLS Libraries}

Users of the XenLinux 2.6 kernel should disable Thread Local Storage
(e.g. by doing a {\tt mv /lib/tls /lib/tls.disabled}) before
attempting to run with a XenLinux kernel.  You can always reenable it
by restoring the directory to its original location (i.e. {\tt mv
  /lib/tls.disabled /lib/tls}).

The TLS implementation uses segmentation in a way that is not
permissable under Xen.  If TLS is not disabled, an emulation mode is
used within Xen which reduces performance substantially and is not
guaranteed to work perfectly.

\section{Test the new install}

It should now be possible to restart the system and use Xen.  Reboot
as usual but choose the new Xen option when the Grub screen appears.

What follows should look much like a conventional Linux boot.  The
first portion of the output comes from Xen itself, supplying low level
information about itself and the machine it is running on.  The
following portion of the output comes from XenLinux itself.

You may see some errors during the XenLinux boot.  These are not
necessarily anything to worry about --- they may result from kernel
configuration differences between your XenLinux kernel and the one you
usually use.

When the boot completes, you should be able to log into your system as
usual.  If you are unable to log in to your system running Xen, you
should still be able to reboot with your normal Linux kernel.


\chapter{Starting a domain}

The first step in creating a new domain is to prepare a root
filesystem for it to boot off.  Typically, this might be stored in a
normal partition, an LVM or other volume manager partition, a disk
file or on an NFS server.

A simple way to do this is simply to boot from your standard OS
install CD and install the distribution into another partition on your
hard drive.

{\em N.b } you can boot with Xen and XenLinux without installing any
special userspace tools but will need to have the prerequisites
described in Section~\ref{sec:prerequisites} and the Xen control tools
installed before you proceed.

\section{From the web interface}

Boot the Xen machine and start Xensv (see Chapter~\ref{cha:xensv} for
more details) using the command: \\
\verb_# xensv start_ \\
This will also start Xend (see Chapter~\ref{cha:xend} for more information).

The domain management interface will then be available at {\tt
http://your\_machine:8080/}.  This provides a user friendly wizard for
starting domains and functions for managing running domains.

\section{From the command line}

Full details of the {\tt xm} tool are found in Chapter~\ref{cha:xm}.

This example explains how to use the \path{xmdefconfig} file.  If you
require a more complex setup, you will want to write a custom
configuration file --- details of the configuration file formats are
included in Chapter~\ref{cha:config}.

The \path{xmexample1} file is a simple template configuration file
for describing a single VM.

The \path{xmexample2} file is a template description that is intended
to be reused for multiple virtual machines.  Setting the value of the
{\tt vmid} variable on the {\tt xm} command line
fills in parts of this template.

Both of them can be found in \path{/etc/xen/}
\subsection{Editing \path{xmdefconfig}}

At minimum, you should edit the following variables in \path{/etc/xen/xmdefconfig}:

\begin{description}
\item[kernel] Set this to the path of the kernel you compiled for use
              with Xen. [e.g. {\tt kernel =
              '/root/xen-2.0.bk/install/boot/vmlinuz-2.4.27-xenU'}]
\item[memory] Set this to the size of the domain's memory in
megabytes. [e.g. {\tt memory = 64 } ]
\item[disk] Set the first entry in this list to calculate the offset
of the domain's root partition, based on the domain ID.  Set the
second to the location of \path{/usr} (if you are sharing it between
domains). [i.e. {\tt disk = ['phy:your\_hard\_drive\%d,sda1,w' \%
(base\_partition\_number + vmid), 'phy:your\_usr\_partition,sda6,r' ]}
\item[dhcp] Uncomment the dhcp variable, so that the domain will
receive its IP address from a DHCP server. [i.e. {\tt dhcp='dhcp'}]
\end{description}

You may also want to edit the {\bf vif} variable in order to choose
the MAC address of the virtual ethernet interface yourself.  For
example: \\ \verb_vif = ['mac=00:06:AA:F6:BB:B3']_\\ If you do not set
this variable, Xend will automatically generate a random MAC address
from an unused range.

If you don't have a \path{xmdefconfig} file, simply create your own 
by copying one of the \path{/etc/xen/xmexample} files.
\subsection{Starting the domain}

The {\tt xm} tool provides a variety of commands for managing domains.
Use the {\tt create} command to start new domains.  To start the
virtual machine with virtual machine ID 1.

\begin{verbatim}
# xm create -c vmid=1
\end{verbatim}

The {\tt -c} switch causes {\tt xm} to turn into the domain's console
after creation.  The {\tt vmid=1} sets the {\tt vmid} variable used in
the {\tt xmdefconfig} file.  The tool uses the
\path{/etc/xen/xmdefconfig} file, since no custom configuration file
was specified on the command line.

\section{Example: ttylinux}

Ttylinux is a very small Linux distribution, designed to
require very few resources.  We will use it as a concrete example of
how to start a Xen domain.  Most users will probably want to install a
more complex mainstream distribution once they have mastered the
basics.

\begin{enumerate}
\item Download the ttylinux disk image from XXX where from?
\item Create a configuration file like the following:
\begin{verbatim}
kernel = "/boot/vmlinuz-2.6.8.1-xenU" # or a 2.4 kernel or a xen0 kernel
memory = 64
name = "ttylinux"
cpu = -1 # leave to Xen to pick
nics = 1
ip = "1.2.3.4"
disk = ['file:/path/to/ttylinux-disk,sda1,w']
root = "/dev/sda1 ro"
\end{verbatim}
\item Now start the domain and connect to its console:
\begin{verbatim}
xm create -f configfile -c
\end{verbatim}
\item Login as root, password root.
\end{enumerate}

\section{Starting / Stopping domains automatically}

It is possible to have certain domains start automatically at boot
time and to have dom0 wait for all running domains to shutdown before
it shuts down the system.

To specify a domain is to start at boot-time, place its
configuration file (or a link to it) under \path{/etc/xen/auto/}.

A Sys-V style init script for RedHat and LSB-compliant systems is
provided and will be automatically copied to \path{/etc/init.d/}
during install.  You can then enable it in the appropriate way for
your distribution.

For instance, on RedHat:

\verb_# chkconfig --add xendomains_

By default, this will start the boot-time domains in runlevels 3, 4
and 5.

You can also use the {\tt service} command to run this script manually, e.g:

\verb_# service xendomains start_

Starts all the domains with config files under /etc/xen/auto/.

\verb_# service xendomains stop_

Shuts down ALL running Xen domains.


\chapter{Domain management tasks}

The previous chapter described a simple example of how to configure
and start a domain.  This chapter summarises the tools available to
manage running domains.

\section{Command line management}

Command line management tasks are also performed using the {\tt xm}
tool.  For online help for the commands available, type:\\
\verb_# xm help_

\subsection{Basic management commands}

The most important {\tt xm} commands are: \\
\verb_# xm list_ : Lists all domains running. \\
\verb_# xm consoles_ : Gives information about the domain consoles. \\
\verb_# xm console_: open a console to a domain.
e.g. \verb_# xm console 1_ (open console to domain 1)

\subsection{\tt xm list}

The output of {\tt xm list} is in rows of the following format:\\
\verb_name domid memory cpu state cputime console_

\begin{description}
\item[name]  The descriptive name of the virtual machine.
\item[domid] The number of the domain ID this virtual machine is running in.
\item[memory] Memory size in megabytes.
\item[cpu]   The CPU this domain is running on.
\item[state] Domain state consists of 5 fields:
  \begin{description}
  \item[r] running
  \item[b] blocked
  \item[p] paused
  \item[s] shutdown
  \item[c] crashed
  \end{description}
\item[cputime] How much CPU time (in seconds) the domain has used so far.
\item[console] TCP port accepting connections to the domain's console.
\end{description}

The {\tt xm list} command also supports a long output format when the
{\tt -l} switch is used.  This outputs the fulls details of the
running domains in Xend's SXP configuration format.

For example, suppose the system is running the ttylinux domain as
described earlier.  The list command should produce output somewhat
like the following:
\begin{verbatim}
# xm list
Name              Id  Mem(MB)  CPU  State  Time(s)  Console
Domain-0           0      251    0  r----    172.2        
ttylinux           5       63    0  -b---      3.0    9605
\end{verbatim}

Here we can see the details for the ttylinux domain, as well as for
domain 0 (which of course is always running).  Note that the console
port for the ttylinux domain is 9605.  This can be connected to by TCP
using a terminal program (e.g. {\tt telnet} or, better, {\tt
xencons}).  The simplest way to connect is to use the {\tt xm console}
command, specifying the domain name or ID.  To connect to the console
of the ttylinux domain, we could use:
\begin{verbatim}
# xm console ttylinux
\end{verbatim}
or:
\begin{verbatim}
# xm console 5
\end{verbatim}

\chapter{Domain filesystem storage}

It is possible to directly export any Linux block device to a virtual,
or to export filesystems / devices to virtual machines using standard
network protocols (e.g. NBD, iSCSI, NFS, etc).  This chapter covers
some of the possibilities.

\section{Warning: Block device sharing}

Block devices should only be shared between domains in a read-only
fashion otherwise the Linux kernels will obviously get very confused
as the file system structure may change underneath them (having the
same partition mounted rw twice is a sure fire way to cause
irreparable damage)!  If you want read-write sharing, export the
directory to other domains via NFS from domain0.

\section{File-backed virtual block devices}

It is possible to use a file in Domain 0 as the primary storage for a
virtual machine.  As well as being convenient, this also has the
advantage that the virtual block device will be {\em sparse} --- space
will only really be allocated as parts of the file are used.  So if a
virtual machine uses only half of its disk space then the file really
takes up half of the size allocated.

For example, to create a 2GB sparse file-backed virtual block device
(actually only consumes 1KB of disk):

\verb_# dd if=/dev/zero of=vm1disk bs=1k seek=2048k count=1_

Make a file system in the disk file: \\
\verb_# mkfs -t ext3 vm1disk_

(when the tool asks for confirmation, answer `y')

Populate the file system e.g. by copying from the current root:
\begin{verbatim}
# mount -o loop vm1disk /mnt
# cp -ax / /mnt
\end{verbatim}
Tailor the file system by editing \path{/etc/fstab},
\path{/etc/hostname}, etc (don't forget to edit the files in the
mounted file system, instead of your domain 0 filesystem, e.g. you
would edit \path{/mnt/etc/fstab} instead of \path{/etc/fstab} ).  For
this example put \path{/dev/sda1} to root in fstab.

Now unmount (this is important!):\\
\verb_# umount /mnt_

In the configuration file set:\\
\verb_disk = ['file:/full/path/to/vm1disk,sda1,w']_

As the virtual machine writes to its `disk', the sparse file will be
filled in and consume more space up to the original 2GB.

\section{NFS Root}

The procedure for using NFS root in a virtual machine is basically the
same as you would follow for a real machine.  NB. the Linux NFS root
implementation is known to have stability problems under high load
(this is not a Xen-specific problem), so this configuration may not be
appropriate for critical servers.

First, populate a root filesystem in a directory on the server machine
--- this can be on another physical machine, or perhaps just another
virtual machine on the same node.

Now, configure the NFS server to export this filesystem over the
network by adding a line to /etc/exports, for instance:

\begin{verbatim}
/export/vm1root      w.x.y.z/m (rw,sync,no_root_squash)
\end{verbatim}

Finally, configure the domain to use NFS root.  In addition to the
normal variables, you should make sure to set the following values in
the domain's configuration file:

\begin{verbatim}
root       = '/dev/nfs'
nfs_server = 'a.b.c.d'       # Substitute the IP for the server here
nfs_root   = '/path/to/root' # Path to root FS on the server machine
\end{verbatim}

The domain will need network access at boot-time, so either statically
configure an IP address (Using the config variables {\tt ip}, {\tt
netmask}, {\tt gateway}, {\tt hostname}) or enable DHCP ({\tt
dhcp='dhcp'}).

\section{LVM-backed virtual block devices}

XXX Put some simple examples here - would be nice if an LVM user could
contribute some, although obviously users would have to read the LVM
docs to do advanced stuff.

\part{Quick Reference}

\chapter{Domain Configuration Files}
\label{cha:config}

XXX Could use a little explanation about possible values

Xen configuration files contain the following standard variables:

\begin{description}
\item[kernel] Path to the kernel image (on the server).
\item[ramdisk] Path to a ramdisk image (optional).
% \item[builder] The name of the domain build function (e.g. {\tt'linux'} or {\tt'netbsd'}.
\item[memory] Memory size in megabytes.
\item[cpu] CPU to assign this domain to.
\item[nics] Number of virtual network interfaces.
\item[vif] List of MAC addresses (random addresses are assigned if not given).
\item[disk] Regions of disk to export to the domain.
\item[dhcp] Set to {\tt 'dhcp'} if you want to DHCP allocate the IP address.
\item[netmask] IP netmask.
\item[gateway] IP address for the gateway (if any).
\item[hostname] Set the hostname for the virtual machine.
\item[root] Set the root device.
\item[nfs\_server] IP address for the NFS server.
\item[nfs\_root] Path of the root filesystem on the NFS server.
\item[extra] Extra string to append to the kernel command line.
\item[restart] Three possible options:
  \begin{description}
  \item[always] Always restart the domain, no matter what
                its exit code is.
  \item[never]  Never restart the domain.
  \item[onreboot] (restart the domain if it requests reboot).
  \end{description}
\end{description}

It is also possible to include Python scripting commands in
configuration files.  This is done in the \path{xmdefconfig} file in
order to handle the {\tt vmid} variable.


\chapter{Xend (Node control daemon)}
\label{cha:xend}

The Xen Daemon (Xend) performs system management functions related to
virtual machines.  It forms a central point of control for a machine
and can be controlled using an HTTP-based protocol.  Xend must be
running in order to start and manage virtual machines.

Xend must be run as root because it needs access to privileged system
management functions.  A small set of commands may be issued on the
Xend command line:

\begin{tabular}{ll}
\verb_# xend start_ & start Xend, if not already running \\
\verb_# xend stop_  & stop Xend if already running       \\
\verb_# xend restart_ & restart Xend if running, otherwise start it \\
\end{tabular}

A SysV init script called {\tt xend} is provided to start Xend at boot
time.  {\tt make install} installs this script in {\path{/etc/init.d}.
To enable it, you have to make symbolic links in the appropriate
runlevel directories or use the {\tt chkconfig} tool, where available.

Once Xend is running, more sophisticated administration can be done
using the xm tool (see Chapter~\ref{cha:xm}) and the experimental
Xensv web interface (see Chapter~\ref{cha:xensv}).

\chapter{The xm tool}
\label{cha:xm}

The xm tool is the primary tool for managing Xen from the console.
The general format of an xm command line is:

\begin{verbatim}
# xm command [switches] [arguments] [variables]
\end{verbatim}

The available {\em switches } and {\em arguments} are dependent on the
{\em command} chosen.  The {\em variables} may be set using
declarations of the form {\tt variable=value} and command line
declarations override any of the values in the configuration file
being used, including the standard variables described above and any
custom variables (for instance, the \path{xmdefconfig} file uses a
{\tt vmid} variable).

The available commands are as follows:

\begin{description}
\item[balloon] Request a domain to adjust its memory footprint.
\item[create] Create a new domain.
\item[destroy] Kill a domain immediately.
\item[list] List running domains.
\item[shutdown] Ask a domain to shutdown.
\item[dmesg] Fetch the Xen (not Linux!) boot output.
\item[consoles] Lists the available consoles.
\item[console] Connect to the console for a domain.
\item[help] Get help on xm commands.
\item[save] Suspend a domain to disk.
\item[restore] Restore a domain from disk.
\item[pause] Pause a domain's execution.
\item[unpause] Unpause a domain.
\item[pincpu] Pin a domain to a CPU.
\item[bvt] Set BVT scheduler parameters for a domain.
\item[bvt\_ctxallow] Set the BVT context switching allowance for the system.
\item[atropos] Set the atropos parameters for a domain.
\item[rrobin] Set the round robin time slice for the system.
\item[info] Get information about the Xen host.
\item[call] Call a Xend HTTP API function directly.
\end{description}

For a detailed overview of switches, arguments and variables to each command
try
\begin{verbatim}
# xm help command
\end{verbatim}

\chapter{Xensv (Web control interface)}
\label{cha:xensv}

Xensv is the experimental web control interface for managing a Xen
machine.  It can be used to perform some (but not yet all) of the
management tasks that can be done using the xm tool.

It can be started using:\\ \verb_# xensv start_ \\ and
stopped using: \verb_# xensv stop_ \\ It will automatically start Xend
if it is not already running.

By default, Xensv will serve out the web interface on port 8080.  This
can be changed by editing {\tt
/usr/lib/python2.3/site-packages/xen/sv/params.py}.

Once Xensv is running, the web interface can be used to manage running
domains and provides a user friendly domain creation wizard.



\chapter{Glossary}

\begin{description}
\item[Atropos]             One of the CPU schedulers provided by Xen.
                           Atropos provides domains with absolute shares
                           of the CPU, with timeliness guarantees and a
                           mechanism for sharing out `slack time'.

\item[BVT]                 The BVT scheduler is used to give proportional
                           fair shares of the CPU to domains.

\item[Exokernel]           A minimal piece of privileged code, similar to
                           a {\bf microkernel} but providing a more
                           `hardware-like' interface to the tasks it
                           manages.  This is similar to a paravirtualising
                           VMM like {\bf Xen} but was designed as a new
                           operating system structure, rather than
                           specifically to run multiple conventional OSs.

\item[Domain]              A domain is the execution context that
                           contains a running { \bf virtual machine }.
                           The relationship between virtual machines
                           and domains on Xen is similar to that between
                           programs and processes in an operating
                           system: a virtual machine is a persistent
                           entity that resides on disk (somewhat like
                           a program).  When it is loaded for execution,
                           it runs in a domain.  Each domain has a
                           { \bf domain ID }.

\item[Domain 0]            The first domain to be started on a Xen
                           machine.  Domain 0 is responsible for managing
                           the system.

\item[Domain ID]           A unique identifier for a { \bf domain },
                           analogous to a process ID in an operating
                           system.  Apart from domain

\item[Full virtualisation] An approach to virtualisation which
                           requires no modifications to the hosted
                           operating system, providing the illusion of
                           a complete system of real hardware devices.

\item[Hypervisor]          An alternative term for { \bf VMM }, used
                           because it means `beyond supervisor',
                           since it is responsible for managing multiple
                           `supervisor' kernels.

\item[Live migration]      A technique for moving a running virtual
                           machine to another physical host, without
                           stopping it or the services running on it.

\item[Microkernel]         A small base of code running at the highest
                           hardware privilege level.  A microkernel is
                           responsible for sharing CPU and memory (and
                           sometimes other devices) between less
                           privileged tasks running on the system.
                           This is similar to a VMM, particularly a
                           {\bf paravirtualising} VMM but typically
                           addressing a different problem space and
                           providing different kind of interface.

\item[NetBSD/Xen]          A port of NetBSD to the Xen architecture.

\item[Paravirtualisation]  An approach to virtualisation which requires
                           modifications to the operating system in
                           order to run in a virtual machine.  Xen
                           uses paravirtualisation but preserves
                           binary compatibility for user space
                           applications.

\item[Shadow pagetables]   A technique for hiding the layout of machine
                           memory from a virtual machine's operating
                           system.  Used in some {\bf VMM}s to provide
                           the illusion of contiguous physical memory,
                           in Xen this is used during
                           {\bf live migration}.

\item[Virtual Machine]     The environment in which a hosted operating
                           system runs, providing the abstraction of a
                           dedicated machine.  A virtual machine may
                           be identical to the underlying hardware (as
                           in { \bf full virtualisation }, or it may
                           differ, as in { \bf paravirtualisation }.

\item[VMM]                 Virtual Machine Monitor - the software that
                           allows multiple virtual machines to be
                           multiplexed on a single physical machine.

\item[Xen]                 Xen is a paravirtualising virtual machine
                           monitor, developed primarily by the
                           Systems Research Group at the University
                           of Cambridge Computer Laboratory.

\item[XenLinux]            Official name for the port of the Linux kernel
                           that runs on Xen.

\end{description}

\part{Advanced Topics}

\chapter{Advanced Network Configuration}

For simple systems with a single ethernet interface with a simple
configuration, the default installation should work `out of the
box'.  More complicated network setups, for instance with multiple
ethernet interfaces and / or existing bridging setups will require
some special configuration.

The purpose of this chapter is to describe the mechanisms provided by
xend to allow a flexible configuration for Xen's virtual networking.

\section{Xen networking scripts}

Xen's virtual networking is configured by 3 shell scripts.  These are
called automatically by Xend when certain events occur, with arguments
to the scripts providing further contextual information.  These
scripts are found by default in \path{/etc/xen}.  The names and
locations of the scripts can be configured in \path{xend-config.sxp}.

\subsection{\path{network}}

This script is called once when Xend is started and once when Xend is
stopped.  Its job is to do any advance preparation required for the
Xen virtual network when Xend starts and to do any corresponding
cleanup when Xend exits.

In the default configuration, this script creates the bridge
`xen-br0' and moves eth0 onto that bridge, modifying the routing
accordingly.

In configurations where the bridge already exists, this script could
be replaced with a link to \path{/bin/true} (for instance).

When Xend exits, this script is called with the {\tt stop} argument,
which causes it to delete the Xen bridge and remove {\tt eth0} from
it, restoring the normal IP and routing configuration.

\subsection{\path{vif-bridge}}

This script is called for every domain virtual interface.  This should
do things like configuring firewalling rules for that interface and
adding it to the appropriate bridge.

By default, this adds and removes VIFs on the default Xen bridge.
This script can be customized to properly deal with more complicated
bridging setups.

\chapter{Advanced Scheduling Configuration}

\section{Scheduler selection}

Xen offers a boot time choice between multiple schedulers.  To select
a scheduler, pass the boot parameter { \tt sched=sched\_name } to Xen,
substituting the appropriate scheduler name.  Details of the schedulers
and their parameters are included below; future versions of the tools
will provide a higher-level interface to these tools.

It is expected that system administrators configure their system to
use the scheduler most appropriate to their needs.  Currently, the BVT
scheduler is the recommended choice, since the Atropos scheduler is
not finished.

\section{Borrowed Virtual Time}

{\tt sched=bvt } (the default) \\ 

BVT provides proportional fair shares of the CPU time.  It has been
observed to penalise domains that block frequently (e.g. IO intensive
domains), but this can be compensated by using warping. 

\subsection{Global Parameters}

\begin{description}
\item[ctx\_allow]
  the context switch allowance is similar to the "quantum"
  in traditional schedulers.  It is the minimum time that
  a scheduled domain will be allowed to run before being
  pre-empted.  This prevents thrashing of the CPU.
\end{description}

\subsection{Per-domain parameters}

\begin{description}
\item[mcuadv]
  the MCU (Minimum Charging Unit) advance determines the
  proportional share of the CPU that a domain receives.  It
  is set inversely proportionally to a domain's sharing weight.
\item[warp]
  the amount of "virtual time" the domain is allowed to warp
  backwards
\item[warpl]
  the warp limit is the maximum time a domain can run warped for
\item[warpu]
  the unwarp requirement is the minimum time a domain must
  run unwarped for before it can warp again
\end{description}

\section{Atropos}

{\tt sched=atropos } \\

Atropos is a Soft Real Time scheduler.  It provides guarantees about
absolute shares of the CPU (with a method for optionally sharing out
slack CPU time on a best-effort basis) and can provide timeliness
guarantees for latency-sensitive domains.

Every domain has an associated period and slice.  The domain should
receive 'slice' nanoseconds every 'period' nanoseconds.  This allows
the administrator to configure both the absolute share of the CPU a
domain receives and the frequency with which it is scheduled.  When
domains unblock, their period is reduced to the value of the latency
hint (the slice is scaled accordingly so that they still get the same
proportion of the CPU).  For each subsequent period, the slice and
period times are doubled until they reach their original values.

Note: don't overcommit the CPU when using Atropos (i.e. don't reserve
more CPU than is available - the utilisation should be kept to
slightly less than 100% in order to ensure predictable behaviour).

\subsection{Per-domain parameters}

\begin{description}
\item[slice]
  The length of time per period that a domain is guaranteed.
\item[period]
  The period over which a domain is guaranteed to receive
  its slice of CPU time.
\item[latency]
  The latency hint is used to control how soon after
  waking up a domain should be scheduled.
\item[xtratime]
  This is a true (1) / false (0) flag that specifies whether
  a domain should be allowed a share of the system slack time.
\end{description}

\section{Round Robin}

{\tt sched=rrobin } \\

The Round Robin scheduler is included as a simple demonstration of
Xen's internal scheduler API.  It is not intended for production use
--- the other schedulers included are all more general and should give
higher throughput.

\subsection{Global parameters}

\begin{description}
\item[rr\_slice]
  The maximum time each domain runs before the next
  scheduling decision is made.
\end{description}

\chapter{Privileged domains}

There are two possible types of privileges:  IO privileges and
administration privileges.

\section{Driver domains (IO Privileges)}

IO privileges can be assigned to allow a domain to drive PCI devices
itself.  This is used to support driver domains.

Setting backend privileges is currently only supported in SXP format
config files (??? is this true - there's nothing in xmdefconfig,
anyhow).  To allow a domain to function as a backend for others,
somewhere within the {\tt vm} element of its configuration file must
be a {\tt backend} element of the form {\tt (backend ({\em type}))}
where {\tt \em type} may be either {\tt netif} or {\tt blkif},
according to the type of virtual device this domain will service.
After this domain has been built, Xend will connect all new and
existing {\em virtual} devices (of the appropriate type) to that
backend.

Note that:
\begin{itemize}
\item a block backend cannot import virtual block devices from other
domains
\item a network backend cannot import virtual network devices from
other domains
\end{itemize}

Thus (particularly in the case of block backends, which cannot import
a virtual block device as their root filesystem), you may need to boot
a backend domain from a ramdisk or a network device.

The privilege to drive PCI devices may also be specified on a
per-device basis.  Xen will assign the minimal set of hardware
privileges to a domain that are required to control its devices.  This
can be configured in either format of configuration file:

\begin{itemize}
\item SXP Format:
  Include {\tt device} elements
  {\tt (device (pci (bus {\em x}) (dev {\em y}) (func {\em z}))) } \\
  inside the top-level {\tt vm} element.  Each one specifies the address
  of a device this domain is allowed to drive ---
  the numbers {\em x},{\em y} and {\em z} may be in either decimal or
  hexadecimal format.
\item Flat Format: Include a list of PCI device addresses of the
  format: \\ {\tt pci = ['x,y,z', ...] } \\ where each element in the
  list is a string specifying the components of the PCI device
  address, separated by commas.  The components ({\tt \em x}, {\tt \em
  y} and {\tt \em z}) of the list may be formatted as either decimal
  or hexadecimal.
\end{itemize}

\section{Administration Domains}

Administration privileges allow a domain to use the `dom0
operations' (so called because they are usually available only to
domain 0).  A privileged domain can build other domains, set scheduling
parameters, etc.

% Support for other administrative domains is not yet available...

\chapter{Debugging}

Xen has a set of debugging features that can be useful to try and
figure out what's going on. Hit 'h' on the serial line (if you
specified a baud rate on the Xen command line) or ScrollLock-h on the
keyboard to get a list of supported commands.

If you have a crash you'll likely get a crash dump containing an EIP
(PC) which, along with an 'objdump -d image', can be useful in
figuring out what's happened.  Debug a Xenlinux image just as you
would any other Linux kernel.

We supply a handy debug terminal program which you can find in
/usr/local/src/xen-2.0.bk/tools/misc/miniterm/
This should be built and executed on another machine that is connected
via a null modem cable. Documentation is included.
Alternatively, if the Xen machine is connected to a serial-port server
then we supply a dumb TCP terminal client, {\tt xencons}.

\chapter{Xen build options}

For most users, the default build of Xen will be adequate.  For some
advanced uses, Xen provides a number of build-time options:

At build time, these options should be set as environment variables or
passed on make's command-line.  For example:

\begin{verbatim}
export option=y; make
option=y make
make option1=y option2=y
\end{verbatim}

\section{List of options}

{\bf verbose=y }\\
Enable debugging messages when Xen detects an unexpected condition.
Also enables console output from all domains. \\
{\bf debug=y }\\
Enable debug assertions.  Implies {\bf verbose=y }.
(Primarily useful for tracing bugs in Xen).        \\
{\bf debugger=y }\\
Enable the in-Xen pervasive debugger (PDB).
This can be used to debug Xen, guest OSes, and
applications. For more information see the 
XenDebugger-HOWTO.                                 \\
{\bf perfc=y }\\
Enable performance-counters for significant events
within Xen. The counts can be reset or displayed
on Xen's console via console control keys.          \\
{\bf trace=y }\\
Enable per-cpu trace buffers which log a range of
events within Xen for collection by control
software.  For more information see the chapter on debugging,
in the Xen Interface Manual.

\chapter{Boot options}

\section{Xen boot options}

These options are used to configure Xen's behaviour at runtime.  They
should be appended to Xen's command line, either manually or by
editing \path{grub.conf}.

{\bf ignorebiostables }\\
 Disable parsing of BIOS-supplied tables. This may help with some
 chipsets that aren't fully supported by Xen. If you specify this
 option then ACPI tables are also ignored, and SMP support is
 disabled. \\

{\bf noreboot } \\
 Don't reboot the machine automatically on errors.  This is
 useful to catch debug output if you aren't catching console messages
 via the serial line. \\

{\bf nosmp } \\
 Disable SMP support.
 This option is implied by 'ignorebiostables'. \\

{\bf noacpi } \\
 Disable ACPI tables, which confuse Xen on some chipsets.
 This option is implied by 'ignorebiostables'. \\

{\bf watchdog } \\
 Enable NMI watchdog which can report certain failures. \\

{\bf noht } \\
 Disable Hyperthreading. \\

{\bf badpage=$<$page number$>$[,$<$page number$>$] } \\
 Specify a list of pages not to be allocated for use 
 because they contain bad bytes. For example, if your
 memory tester says that byte 0x12345678 is bad, you would
 place 'badpage=0x12345' on Xen's command line (i.e., the
 last three digits of the byte address are not
 included!). \\

{\bf com1=$<$baud$>$,DPS[,$<$io\_base$>$,$<$irq$>$] \\
 com2=$<$baud$>$,DPS[,$<$io\_base$>$,$<$irq$>$] } \\
 Xen supports up to two 16550-compatible serial ports.
 For example: 'com1=9600,8n1,0x408,5' maps COM1 to a
 9600-baud port, 8 data bits, no parity, 1 stop bit,
 I/O port base 0x408, IRQ 5.
 If the I/O base and IRQ are standard (com1:0x3f8,4;
 com2:0x2f8,3) then they need not be specified. \\

{\bf console=$<$specifier list$>$ } \\
 Specify the destination for Xen console I/O.
 This is a comma-separated list of, for example:
\begin{description}
 \item[vga]  use VGA console and allow keyboard input
 \item[com1] use serial port com1
 \item[com2H] use serial port com2. Transmitted chars will
   have the MSB set. Received chars must have
   MSB set.
 \item[com2L] use serial port com2. Transmitted chars will
   have the MSB cleared. Received chars must
   have MSB cleared.
\end{description}
 The latter two examples allow a single port to be
 shared by two subsystems (e.g. console and
 debugger). Sharing is controlled by MSB of each
 transmitted/received character.
 [NB. Default for this option is 'com1,tty'] \\

{\bf conswitch=$<$switch-char$><$auto-switch-char$>$ } \\
 Specify how to switch serial-console input between
 Xen and DOM0. The required sequence is CTRL-<switch-char>
 pressed three times. Specifying '`' disables switching.
 The <auto-switch-char> specifies whether Xen should
 auto-switch input to DOM0 when it boots -- if it is 'x'
 then auto-switching is disabled.  Any other value, or
 omitting the character, enables auto-switching.
 [NB. Default for this option is 'a'] \\

{\bf nmi=xxx } \\
 Specify what to do with an NMI parity or I/O error. \\
 'nmi=fatal':  Xen prints a diagnostic and then hangs. \\
 'nmi=dom0':   Inform DOM0 of the NMI. \\
 'nmi=ignore': Ignore the NMI. \\

{\bf dom0\_mem=xxx } \\
 Set the maximum amount of memory for domain0. \\

{\bf tbuf\_size=xxx } \\
 Set the size of the per-cpu trace buffers, in pages
 (default 1).  Note that the trace buffers are only
 enabled in debug builds.  Most users can ignore
 this feature completely. \\

{\bf sched=xxx } \\
 Select the CPU scheduler Xen should use.  The current
 possibilities are 'bvt', 'atropos' and 'rrobin'.  The
 default is 'bvt'.  For more information see
 Sched-HOWTO.txt. \\

{\bf pci\_dom0\_hide=(xx.xx.x)(yy.yy.y)... } \\
Hide selected PCI devices from domain 0 (for instance, to stop it
taking ownership of them so that they can be driven by another
domain).  Device IDs should be given in hex format.  Bridge devices do
not need to be hidden --- they are hidden implicitly, since guest OSes
do not need to configure them.

\section{XenLinux Options}

{\bf xencons=xxx}
Specify the device node to
which the Xen virtual console driver is attached: \\
 'xencons=off': disable virtual console \\
 'xencons=tty': attach console to /dev/tty1 (tty0 at boot-time) \\
 'xencons=ttyS': attach console to /dev/ttyS0\\
The default is ttyS for dom0 and tty for all other domains.

\chapter{Further Support}

If you have questions that are not answered by this manual, the
sources of information listed below may be of interest to you.  Note
that bug reports, suggestions and contributions related to the
software (or the documentation) should be sent to the Xen developers'
mailing list (address below).

\section{Other documentation}

For developers interested in porting operating systems to Xen, the
{\em Xen Interface Manual} is distributed in the \path{docs/}
directory of the Xen source distribution.  Various HOWTOs are
available in \path{docs/HOWTOS} but this content is being integrated
into this manual.

\section{Online references}

The official Xen web site is found at: \\
{\tt
http://www.cl.cam.ac.uk/Research/SRG/netos/xen/] }.

Links to other
documentation sources are listed at: \\ {\tt
http://www.cl.cam.ac.uk/Research/SRG/netos/xen/documentation.html}.

\section{Mailing lists}

There are currently two official Xen mailing lists:

\begin{description}
\item[xen-devel@lists.sourceforge.net] Used for development
discussions and requests for help.  Subscribe at: \\
{\tt http://lists.sourceforge.net/mailman/listinfo/xen-devel}
\item[xen-announce@lists.sourceforge.net] Used for announcements only.
Subscribe at: \\
{\tt http://lists.sourceforge.net/mailman/listinfo/xen-announce}
\item[xen-changelog@lists.sourceforge.net]  Changelog feed
from the unstable and 2.0 trees - developer oriented.  Subscribe at: \\
{\tt http://lists.sourceforge.net/mailman/listinfo/xen-changelog}
\end{description}

Although there is no specific user support list, the developers try to
assist users who post on xen-devel.  As the bulk of traffic on this
list increases, a dedicated user support list may be introduced.

\appendix

\chapter{Installing Debian}

The Debian project provides a tool called {\tt debootstrap} which
allows a base Debian system to be installed into a filesystem without
requiring the host system to have any Debian-specific software (such
as {\tt apt}).

Here's some info how to install Debian 3.1 (Sarge) for an unprivileged
Xen domain:

\begin{enumerate}
\item Set up Xen 2.0 and test that it's working, as described earlier in
      this manual.

\item Create disk images for root-fs and swap (alternatively, you
      might create dedicated partitions, LVM logical volumes, etc. if
      that suits your setup).
\begin{verbatim}  
dd if=/dev/zero of=/path/diskimage bs=1024k count=size_in_mbytes
dd if=/dev/zero of=/path/swapimage bs=1024k count=size_in_mbytes
\end{verbatim}
      If you're going to use this filesystem / diskimage only as a
      `template' for other vm diskimages, something like 300 MB should
      be enough.. (of course it depends what kind of packages you are
      planning to install to the template)

\item Create the filesystem and initialise the swap image
\begin{verbatim}
mkfs.ext3 /path/diskimage
mkswap /path/swapimage
\end{verbatim}

\item Mount the diskimage for installation
\begin{verbatim}
mount -o loop /path/diskimage /mnt/disk
\end{verbatim}

\item Install {\tt debootstrap}

Make sure you have debootstrap installed on the host.  If you are
running Debian sarge (3.1 / testing) or unstable you can install it by
running {\tt apt-get install debootstrap}.  Otherwise, it can be
downloaded from the Debian project website.

\item Install debian base to the diskimage:
\begin{verbatim}
debootstrap --arch i386 sarge /mnt/disk  \
            http://ftp.<countrycode>.debian.org/debian
\end{verbatim}

You can use any other Debian http/ftp mirror you want.

\item When debootstrap completes successfully, modify settings:
\begin{verbatim}
chroot /mnt/disk /bin/bash
\end{verbatim}

Edit the following files using vi or nano and make needed changes:
\begin{verbatim}
/etc/hostname
/etc/hosts
/etc/resolv.conf
/etc/network/interfaces
/etc/networks
\end{verbatim}

Set up access to the services, edit:
\begin{verbatim}
/etc/hosts.deny
/etc/hosts.allow
/etc/inetd.conf
\end{verbatim}

Add Debian mirror to:   
\begin{verbatim}
/etc/apt/sources.list
\end{verbatim}

Create fstab like this:
\begin{verbatim}
/dev/sda1       /       ext3    errors=remount-ro       0       1
/dev/sda2       none    swap    sw                      0       0
proc            /proc   proc    defaults                0       0
\end{verbatim}

Logout

\item      Umount the diskimage
\begin{verbatim}
umount /mnt/disk
\end{verbatim}

\item Create Xen 2.0 configuration file for the new domain. You can
        use the example-configurations coming with xen as a template.

        Make sure you have the following set up:
\begin{verbatim}
disk = [ 'file:/path/diskimage,sda1,w', 'file:/path/swapimage,sda2,w' ]
root = "/dev/sda1 ro"
\end{verbatim}

\item Start the new domain
\begin{verbatim}
xm create -f domain_config_file
\end{verbatim}

Check that the new domain is running:
\begin{verbatim}
xm list
\end{verbatim}

\item   Attach to the console of the new domain.
        You should see something like this when starting the new domain:

\begin{verbatim}
Started domain testdomain2, console on port 9626
\end{verbatim}
        
        There you can see the ID of the console: 26. You can also list
        the consoles with {\tt xm consoles"}. (ID is the last two
        digits of the portnumber.)

        Attach to the console:

\begin{verbatim}
xm console 26
\end{verbatim}

        or by telnetting to the port 9626 of localhost (the xm console
        progam works better).

\item   Log in and run base-config

        As a default there's no password for the root.

        Check that everything looks OK, and the system started without
        errors.  Check that the swap is active, and the network settings are
        correct.

        Run {\tt /usr/sbin/base-config} to set up the Debian settings.

        Set up the password for root using passwd.

\item     Done. You can exit the console by pressing {\tt Ctrl + ]}

\end{enumerate}

If you need to create new domains, you can just copy the contents of
the `template'-image to the new disk images, either by mounting the
template and the new image, and using {\tt cp -a} or {\tt tar} or by
simply copying the image file.  Once this is done, modify the
image-specific settings (hostname, network settings, etc).

\chapter{Installing Xen / XenLinux on Redhat / Fedora}

When using Xen / Xenlinux on a standard Linux distribution there are
a couple of things to watch out for:

Note that, because domains>0 don't have any privileged access at all,
certain commands in the default boot sequence will fail e.g. attempts
to update the hwclock, change the console font, update the keytable
map, start apmd (power management), or gpm (mouse cursor).  Either
ignore the errors (they should be harmless), or remove them from the
startup scripts.  Deleting the following links are a good start:
S24pcmcia S09isdn S17keytable S26apmd S85gpm.

If you want to use a single root file system that works cleanly for
domain0 and domains>0, a useful trick is to use different 'init' run
levels. For example, on the Xen Demo CD we use run level 3 for domain
0, and run level 4 for domains>0. This enables different startup
scripts to be run in depending on the run level number passed on the
kernel command line.

If you're going to use NFS root files systems mounted either from an
external server or from domain0 there are a couple of other gotchas.
The default /etc/sysconfig/iptables rules block NFS, so part way
through the boot sequence things will suddenly go dead.

If you're planning on having a separate NFS /usr partition, the RH9
boot scripts don't make life easy - they attempt to mount NFS file
systems way to late in the boot process. The easiest way I found to do
this was to have a '/linuxrc' script run ahead of /sbin/init that
mounts /usr:

\begin{verbatim}
 #!/bin/bash
 /sbin/ipconfig lo 127.0.0.1
 /sbin/portmap
 /bin/mount /usr
 exec /sbin/init "$@" <>/dev/console 2>&1
\end{verbatim}

The one slight complication with the above is that /sbib/portmap is
dynamically linked against /usr/lib/libwrap.so.0 Since this is in
/usr, it won't work. This can be solved by copying the file (and link)
below the /usr mount point, and just let the file be 'covered' when
the mount happens.

In some installations, where a shared read-only /usr is being used, it
may be desirable to move other large directories over into the
read-only /usr. For example, you might replace /bin /lib and /sbin
with links into /usr/root/bin /usr/root/lib and /usr/root/sbin
respectively. This creates other problems for running the /linuxrc
script, requiring bash, portmap, mount, ifconfig, and a handful of
other shared libraries to be copied below the mount point - little
statically linked C program would solve this problem.


\end{document}


%% Other stuff without a home

%% Instructions Re Python API

%% Other Control Tasks using Python
%% ================================

%% A Python module 'Xc' is installed as part of the tools-install
%% process. This can be imported, and an 'xc object' instantiated, to
%% provide access to privileged command operations:

%% # import Xc
%% # xc = Xc.new()
%% # dir(xc)
%% # help(xc.domain_create)

%% In this way you can see that the class 'xc' contains useful
%% documentation for you to consult.

%% A further package of useful routines (xenctl) is also installed:

%% # import xenctl.utils
%% # help(xenctl.utils)

%% You can use these modules to write your own custom scripts or you can
%% customise the scripts supplied in the Xen distribution.



% Explain about AGP GART


%% If you're not intending to configure the new domain with an IP address
%% on your LAN, then you'll probably want to use NAT. The
%% 'xen_nat_enable' installs a few useful iptables rules into domain0 to
%% enable NAT. [NB: We plan to support RSIP in future]




%% Installing the file systems from the CD
%% =======================================

%% If you haven't got an existing Linux installation onto which you can
%% just drop down the Xen and Xenlinux images, then the file systems on
%% the CD provide a quick way of doing an install. However, you would be
%% better off in the long run doing a proper install of your preferred
%% distro and installing Xen onto that, rather than just doing the hack
%% described below:

%% Choose one or two partitions, depending on whether you want a separate
%% /usr or not. Make file systems on it/them e.g.: 
%%   mkfs -t ext3 /dev/hda3
%%   [or mkfs -t ext2 /dev/hda3 && tune2fs -j /dev/hda3 if using an old
%% version of mkfs]

%% Next, mount the file system(s) e.g.:
%%   mkdir /mnt/root && mount /dev/hda3 /mnt/root
%%   [mkdir /mnt/usr && mount /dev/hda4 /mnt/usr]
  
%% To install the root file system, simply untar /usr/XenDemoCD/root.tar.gz:
%%   cd /mnt/root && tar -zxpf /usr/XenDemoCD/root.tar.gz

%% You'll need to edit /mnt/root/etc/fstab to reflect your file system
%% configuration. Changing the password file (etc/shadow) is probably a
%% good idea too.

%% To install the usr file system, copy the file system from CD on /usr,
%% though leaving out the "XenDemoCD" and "boot" directories:
%%   cd /usr && cp -a X11R6 etc java libexec root src bin dict kerberos local sbin tmp doc include lib man share /mnt/usr

%% If you intend to boot off these file systems (i.e. use them for
%% domain 0), then you probably want to copy the /usr/boot directory on
%% the cd over the top of the current symlink to /boot on your root
%% filesystem (after deleting the current symlink) i.e.:
%%   cd /mnt/root ; rm boot ; cp -a /usr/boot .
